\chapter{Stage 2: The Episodic Memory Store}
\label{ch:memory}

Stage 1 decided how ready the model is. Stage 2 asks a different question: have we
answered something like this well before? The episodic memory store is \caem{}'s
answer to the first structural flaw of \Cref{ch:problem}, statelessness. It is the
component that lets the system say ``I have seen this, and last time I got it
right,'' and it is the substrate that every later stage either reads from or writes
to. This chapter covers the store as a whole: what a single memory holds, how a
query finds its nearest match, how a new memory is admitted, and how the store
keeps itself small and trustworthy over time.

\begin{intuition}[A notebook of verified answers]
Picture a researcher who keeps a notebook of every question they have settled, with
the reasoning that settled it and a note of how sure they were. When a new question
arrives, they first flip to the most similar past entry. If it is close enough and
was confidently resolved, they reuse it instead of redoing the work. The notebook is
curated: near-duplicate entries are merged, shaky entries are re-checked when the
researcher learns more, and if it ever overflows, the least useful pages are torn
out first. The episodic memory store is that notebook, made precise enough for a
machine to keep.
\end{intuition}

\section{What a single memory holds}
\label{sec:memory-episode}

The unit of memory is an \emph{episode}: a verified triple of question, reasoning,
and answer, wrapped in the metadata that lets the rest of the system reason about
it. The fields fall into three groups, and the grouping is not cosmetic, because the
groups have different mutability rules.

\begin{table}[htbp]
\centering
\caption[Anatomy of an episode]{The fields of a stored episode, by group. Content
is frozen at storage time because the fine-tune of \Cref{ch:cycle} trains directly
on it; quality and usage evolve across cycles.}
\label{tab:episode}
\small
\renewcommand{\arraystretch}{1.22}
\begin{tabularx}{\textwidth}{@{}l >{\raggedright\arraybackslash}X@{}}
\toprule
\textbf{Group} & \textbf{Fields and meaning} \\
\midrule
\textbf{Content} (immutable)
  & The \emph{question}, the \emph{reasoning chain} that produced the answer, the
    short \emph{answer} string, and the $768$-dimensional unit-normalised
    \emph{embedding}. Plus provenance: the \emph{storage cycle}, a timestamp, the
    \emph{source benchmark}, and any benchmark tags merged in during consolidation.
    These never change after storage, because the self-improvement loop trains on
    them and a moving target would corrupt the dataset. \\
\addlinespace[2pt]
\textbf{Quality} (mutable)
  & The stored trust score $\ustored$, the verifier signals recorded on the episode
    ($\utoken$, $\udrop$, $\uint$, $\savg$, $\pentail$, $\pgmax$, $\pgmean$,
    $\pgatom$, $\qarel$) together with the retired-but-recorded $h_{\text{norm}}$
    diagnostic, the gate
    \emph{decision} (one of store, defer, abstain, discard), and a confabulation
    flag. These are overwritten when the episode is re-scored at a cycle boundary. \\
\addlinespace[2pt]
\textbf{Usage} (mutable)
  & The \emph{retrieval count}, the running \emph{success rate} of times the episode
    was served and accepted, a \emph{hit counter} since its last re-check, and a
    \emph{retroverified} flag. These track how the episode performs in service and
    drive both the trust-feedback update and the pruning score. \\
\bottomrule
\end{tabularx}
\end{table}

\begin{intuition}[Why store the reasoning, not just the answer]
The most valuable field is often the one a cache would omit: the reasoning chain. A
plain lookup table would store only the question and the answer. \caem{} stores the
chain of reasoning too, because that is the transferable knowledge the
self-improvement loop of \Cref{ch:cycle} learns from. An answer tells the model
\emph{what}; the chain tells it \emph{how}, and the ``how'' is what generalises to
the next, slightly different question. The memory is therefore not just a cache of
answers but a corpus of worked solutions.
\end{intuition}

The trust score $\ustored$ deserves a forward pointer rather than a definition here.
It is the calibrated composite produced by the verifier and gate of
\Cref{ch:verifier,ch:composite}; for this chapter it is enough to read it as a number
in $[0,1]$ that says how much the system trusts this episode's answer. Stage 2 only
\emph{consumes} it, during routing; it is \emph{produced} downstream.

\section{The per-query operation: encode and probe}
\label{sec:memory-probe}

In the per-query pipeline, Stage 2's active job is small and fast. The query is
encoded into its $768$-dimensional embedding by the sentence encoder of
\Cref{sec:prereq-embeddings}, and that vector is used to probe the store for its
single nearest neighbour. Because every vector is unit-normalised, the
nearest-neighbour search by inner product returns exactly the highest cosine
similarity, and the search returns both the matched episode and that similarity
score.

\begin{lstlisting}[caption={The nearest-neighbour probe, condensed from \texttt{caem/memory/store.py}.}, label={lst:search}]
def search(self, embedding, k=1):
    if self.is_empty:
        return []                                # cold start: nothing to match
    vec = self._to_faiss_matrix(embedding)       # (1, 768) float32, unit-norm
    sims, ids = self._index.search(vec, k)       # inner product == cosine
    return [(self._metadata[int(i)], float(s))   # (episode, similarity)
            for s, i in zip(sims[0], ids[0]) if i != -1]
\end{lstlisting}

Two things are worth noticing. First, an empty store returns nothing, which is the
cold-start case: on the very first cycle the memory is empty, so every query falls
through to generation, exactly as we saw for the FEVER claim in \Cref{ch:bigpicture}.
Second, the search returns the episode object itself, not just an identifier, so the
router can immediately read the neighbour's stored trust score without a second
lookup. A companion form of the call also returns the integer identifier, which the
pipeline uses when it later needs to update that episode's usage statistics after a
Tier 1 hit.

\begin{fromcode}[The substrate underneath]
The search rides on the FAISS index described in \Cref{sec:prereq-faiss}. The store
bootstraps as an \emph{exact} inner-product index so that early cycles, when the
store is small, get exact nearest neighbours, and it \emph{auto-promotes} to the
approximate inverted-file product-quantised index once it holds enough vectors to
train the clustering. The index is wrapped so that episodes can be deleted by
identifier, which matters for the pruning and consolidation operations later in this
chapter. None of this is visible to the caller: the probe returns the nearest
episode and a cosine score regardless of which index is active underneath.
\end{fromcode}

\section{Admission: the novelty gate}
\label{sec:memory-novelty}

When the gate of \Cref{ch:composite} decides an answer is worth storing, the store
does not admit it blindly. It first checks \emph{novelty}: if the candidate's nearest
existing neighbour is too similar, the candidate is judged a near-duplicate and is not
stored. ``Too similar'' is a cosine similarity above $0.95$.

\begin{intuition}[Why duplicates are actively harmful]
A cache would happily store the same answer twice; it costs only space. For \caem{}
a duplicate costs more than space, because the memory is also the training set. Storing
the same verified answer many times would let a handful of popular questions dominate
the fine-tune, teaching the model their phrasing at the expense of breadth. The novelty
gate keeps the memory a set of \emph{distinct} solved problems rather than a frequency
table of the common ones, which is what keeps the downstream training signal balanced.
\end{intuition}

The novelty check reuses the very same nearest-neighbour search as the per-query probe:
it asks for the top neighbour and rejects the candidate if that neighbour's similarity
exceeds the threshold. The threshold sits high, at $0.95$, deliberately. Paraphrases
that are merely related should still be stored as separate episodes, because they may
carry different reasoning; only near-identical restatements are screened out at
admission. Genuine paraphrase clusters are handled later and more carefully by
consolidation (\Cref{sec:memory-cycle-ops}), which can inspect the answers before
merging rather than rejecting at the door.

\section{Capacity and pruning}
\label{sec:memory-pruning}

The store has a hard capacity of one million episodes. Long before that, when it
reaches ninety-five percent full, it triggers a prune that removes the least valuable
fifth of its episodes. ``Least valuable'' is not ``oldest'' or ``least confident''
alone; it is a composite \emph{Value} score that balances how good an episode is
against how useful it has proven and how recent it is.

\begin{precise}[The Value score]
Each episode's Value combines an importance term and a recency term,
\[
  \text{Value} \;=\; 0.70\cdot\text{Importance} \;+\; 0.30\cdot\text{Recency},
\]
where importance is a weighted blend of four cues,
\[
  \text{Importance}
  \;=\;
  0.40\,\varphi \;+\; 0.30\,\frac{r}{r_{\max}} \;+\; 0.20\,\text{success} \;+\; 0.10\,\ustored,
\]
and recency decays exponentially with the episode's age in seconds,
\[
  \text{Recency} \;=\; \exp(-0.01\cdot \text{age}).
\]
Here $\varphi = (\text{storage cycle} + 1)/(\text{number of cycles} + 1)$ is a proxy
for how recently the episode was stored, $r/r_{\max}$ is the episode's retrieval count
normalised by the busiest episode's, $\text{success}$ is its running acceptance rate
in service, and $\ustored$ is its trust score. When a prune fires, episodes are sorted
by Value and the bottom twenty percent are removed.
\end{precise}

\begin{workedexample}[Which episode survives a prune]
As an illustration (the capacity prune was never reached in the realised run, whose
store peaked near ten thousand episodes against a one-million capacity),
consider two episodes competing to stay. Episode A was stored in an early cycle, has
never been retrieved, and carries a middling trust score. Its $\varphi$ is small, its
retrieval and success terms are zero, and its recency has decayed, so its Value is
low and it is a prime candidate for removal. Episode B was stored recently, has been
retrieved many times and accepted almost every time, and carries a high trust score.
Every term in its importance is large and its recency is near one, so its Value is
high and it is retained. The lesson is that the store does not just keep its most
confident episodes; it keeps the ones that have \emph{earned their place} by being
useful in service, which is a stronger signal of value than confidence at storage
time alone.
\end{workedexample}

\section{Living metadata: trust that adapts to service}
\label{sec:memory-feedback}

An episode's trust score is not frozen the moment it is stored. Every time the episode
is served from memory and the answer is accepted or overridden, the store updates the
episode's usage statistics, and once the episode has been retrieved enough times, it
also nudges the trust score itself.

\begin{lstlisting}[caption={Retrieval feedback, condensed from \texttt{caem/memory/store.py}.}, label={lst:feedback}]
entry.retrieval_count += 1
entry.success_rate = (old_rate * n + float(was_accepted)) / (n + 1)   # running mean

if entry.retrieval_count >= 5:                  # feedback_loop_min_retrievals
    eta = 0.01                                  # feedback_loop_eta
    if was_accepted:
        entry.u_stored += eta * (1.0 - entry.u_stored)    # nudge toward 1
    else:
        entry.u_stored -= eta * entry.u_stored            # nudge toward 0
\end{lstlisting}

The success rate is a running mean, so it reflects the episode's whole service history,
not just its last use. The trust nudge is deliberately small and deliberately gated: it
applies only after at least five retrievals, so that a single acceptance or rejection
cannot swing an episode's standing, and even then it moves the score by only a small
fraction of the remaining distance to one or zero. The effect accumulates: an episode
that is repeatedly served and accepted drifts slowly upward in trust, while one that is
repeatedly overridden drifts down, eventually toward the pruning threshold. There is
also a separate hit counter that, once an episode has been served ten times since its
last re-check, flags it for forced re-verification at the next cycle boundary, on the
principle that a popular episode that is quietly wrong does the most damage and is worth
re-checking early.

\section{What the store does at a cycle boundary}
\label{sec:memory-cycle-ops}

Two of the store's most important operations are not driven by the per-query pipeline at
all. They are invoked by the self-improvement loop of \Cref{ch:cycle} once per cycle,
and we introduce them here only because they are the store's own methods; their full
treatment belongs to that chapter.

\begin{description}[leftmargin=0pt, labelsep=0.5em, style=nextline]
  \item[Retroactive re-verification.] Every stored episode is re-scored through the
  recalibrated verifier under the freshly fine-tuned model. An episode whose new trust
  score falls below a floor is pruned; otherwise its quality fields are overwritten with
  the new reading, which may raise or lower its trust. This is how the memory benefits
  from the model getting better: yesterday's borderline episode can be re-judged
  confidently today, and a stale confident episode can be caught and demoted. In the
  deployed run this pass pruned a declining fraction at each committed cycle, from about
  thirteen percent at the first cycle to about four percent by the last (about $13$, $9$,
  $7$, $4$ percent), the signature of a store whose composition shifts toward
  higher-confidence survivors as the model improves; the prune floor of $0.50$ sits below
  the storage cut-point, so re-verification removes only entries that have genuinely lost
  confidence.
  \item[Consolidation.] Near-paraphrase episodes are clustered, and each cluster is
  collapsed onto a single representative, the one with the highest trust. The merge is
  guarded: it proceeds only if the cluster's answers agree after normalisation and their
  trust scores do not spread too widely, and it never merges across the
  training-versus-transfer boundary, so that consolidation cannot leak held-out content
  into the training pool. The representative inherits the cluster's combined retrieval
  history and benchmark tags. In the deployed run the clustering threshold sat at cosine
  $0.92$ over the nearest twenty neighbours; the answer-equality guard and a trust-spread
  guard, which rejected any cluster whose trust range exceeded $0.15$, screened most
  candidates, and the cross-pool guard rejected every candidate straddling the
  training-versus-transfer boundary. Across the trajectory the pass merged $57$ clusters
  and removed $58$ episodes. Note that the consolidation threshold of $0.92$ stays
  distinct from the admission novelty gate of $0.95$ described in
  \Cref{sec:memory-novelty}: consolidation inspects answers and trust before merging,
  whereas the novelty gate rejects only near-identical restatements at the door.
\end{description}

\begin{figure}[htbp]
\centering
\begin{tikzpicture}[
  font=\small\sffamily,
  ep/.style={rectangle, rounded corners=5pt, draw=cbExample!70!black,
    fill=cbExample!12, line width=1.1pt, align=center, inner sep=6pt,
    text width=2.7cm, minimum height=1.1cm, drop shadow={shadow xshift=0.3ex,
    shadow yshift=-0.3ex, fill=black!10}},
  act/.style={rectangle, rounded corners=3pt, draw=cbIntuition!75!black,
    fill=cbIntuition!10, line width=0.9pt, align=center, inner sep=4pt,
    text width=2.5cm, font=\footnotesize, minimum height=0.85cm},
  fin/.style={rectangle, rounded corners=2pt, draw=cbInk!55, fill=cbInk!6,
    line width=0.9pt, align=center, inner sep=4pt, font=\footnotesize},
  flow/.style={-{Stealth[length=2.4mm]}, line width=1.1pt, draw=cbInk!75},
  gone/.style={-{Stealth[length=2.4mm]}, line width=1.1pt, draw=cbPitfall!75, dashed},
]
  \node[fin] (gate) at (-4.2,0) {gate says\\\textbf{store}};
  \node[ep]  (live) at (0.4,0) {\textbf{Live episode}\\in the store};
  % feedback self-loop
  \node[act] (fb) at (0.4,2.3) {served $\to$ trust nudge,\\usage stats};
  % cycle-boundary ops
  \node[act] (retro) at (4.8,1.4) {cycle boundary:\\re-verify};
  \node[act] (cons)  at (4.8,-0.0) {cycle boundary:\\consolidate};
  \node[act] (prune) at (4.8,-1.4) {capacity 95\%:\\prune};
  % outcomes
  \node[fin] (rep)  at (9.0,0.7) {merged into\\representative};
  \node[fin, draw=cbPitfall!60, text=cbPitfall!70!black] (rm) at (9.0,-0.9) {removed};

  \draw[flow] (gate) -- node[above,font=\scriptsize,text=cbInk!60]{novelty $\le 0.95$} (live);
  \draw[flow] (live.30) to[bend right=25] (fb.330);
  \draw[flow] (fb.210) to[bend right=25] (live.150);
  \draw[flow] (live) |- (retro);
  \draw[flow] (live) -- (cons);
  \draw[flow] (live) |- (prune);
  \draw[flow] (retro) -| node[pos=0.25,above,font=\scriptsize,text=cbExample!55!black]{pass} (rep.north west);
  \draw[gone] (retro) -- node[pos=0.2,above right,font=\scriptsize,text=cbPitfall!70!black]{below floor} (rm);
  \draw[flow] (cons) -- (rep);
  \draw[gone] (cons) -- node[pos=0.2,below right,font=\scriptsize,text=cbPitfall!70!black]{non-rep} (rm);
  \draw[gone] (prune) -- (rm);
\end{tikzpicture}
\caption[The life of an episode]{The life of an episode. It is admitted only if it
clears the novelty gate, then lives in the store being served (each serve nudges its
trust and usage stats). At each cycle boundary it is re-verified (kept with an updated
trust, or removed if it falls below the floor) and may be consolidated into a
representative; under capacity pressure the least valuable episodes are pruned. The
dashed paths are the only ways an episode leaves the store. (Schematic.)}
\label{fig:episode-life}
\end{figure}

\begin{pitfall}[The memory is curated, not merely accumulated]
It would be a mistake to picture the store as an append-only log that only ever grows.
It is continuously curated. New episodes are screened at admission, served episodes have
their trust nudged by experience, the whole store is re-judged each cycle as the model
improves, paraphrase clusters are merged, and the least valuable episodes are pruned
under capacity pressure. The store's quality is maintained by these forces acting
together, which is what lets a single fixed trust threshold keep meaning the same thing
as the store ages.
\end{pitfall}

\section{The store survives a rollback}
\label{sec:memory-asymmetry}

One property from \Cref{ch:bigpicture} is worth restating in the store's own terms,
because it is the reason the memory is the system's most durable asset. When a cycle's
fine-tune fails the retention guard and the model weights are rolled back, the episodic
memory is \emph{not} rolled back with them. Every episode admitted during that cycle
survives. This fired once in the deployed run; at one cycle the worst of the three
retention probes (the held-out open-text factoid probe) fell to a ratio of $0.886$
against its pristine baseline, below the fixed floor of $0.93$, while the
general-knowledge probe held at $1.018$ and the commonsense probe at $0.982$; the
adapter was restored to the previous cycle's weights, yet the store was untouched and
kept growing across the boundary, from about five thousand to nearly ten thousand
episodes by the final cycle.

\begin{precise}[Why memory growth is safe to keep even on a failed cycle]
The asymmetry is principled. Each episode in the store passed the four-outcome gate at
storage time, so memory growth is monotone and high-precision by construction: adding a
gated episode can only add a trustworthy fact. A weight update has no such guarantee; it
is a global change that can erode unrelated capabilities, which is the catastrophic
forgetting the whole loop is built to avoid. So the system protects the verified,
accumulated asset (the memory) and is willing to discard the risky, reversible one (the
adapter). This is why a rollback costs a cycle's training but never a cycle's
remembered knowledge. The store is also checkpointed to disk each cycle, as a paired
index and metadata file, so the accumulated memory is recoverable independently of the
model.
\end{precise}

\begin{bythenumbers}[Episodic memory constants]
\begin{itemize}[leftmargin=1.3em, itemsep=1pt]
  \item Embedding: $768$-dimensional, unit-normalised; search by cosine (inner product).
  \item Capacity: $1{,}000{,}000$ episodes; prune trigger at $95\%$ full; prune the bottom $20\%$ by Value.
  \item Novelty: candidates with a nearest neighbour above cosine $0.95$ are not stored.
  \item Value $= 0.70\cdot$Importance$+0.30\cdot$Recency; recency decay rate $0.01$ per second of age.
  \item Trust feedback: nudge step $0.01$, applied only after $\geq 5$ retrievals.
  \item Forced re-check: flagged after $10$ Tier 1 serves since the last re-verification.
\end{itemize}
\end{bythenumbers}

\begin{defence}[If an examiner asks: ``how is this different from a retrieval cache?'']
A cache stores answers and serves them back; the episodic store does three things a
cache does not. It stores the \emph{reasoning}, not just the answer, because the memory
doubles as the training corpus for the self-improvement loop. It carries a
\emph{calibrated trust score} on every entry, so the router can weigh a match by quality
and not only by similarity. And it is \emph{curated against its own model}: each cycle
re-judges, merges, and prunes its contents as the model improves, so the store's
precision is maintained rather than left to decay. A cache is a passive lookup table;
this is an actively maintained, quality-scored knowledge base that happens to also
answer lookups.
\end{defence}

The store can now tell the router, for any query, both how similar the nearest episode
is and how much that episode is trusted. \Cref{ch:router} takes up Stage 3, where those
two numbers are fused into a single dispatch score and a tier is finally chosen.
